\chapter{树}

\section{树的基础}

\subsection{树的重心、中心与直径}

该封装用于求树的重心、中心和直径。重心由各连通块大小决定，中心由到最远节点的距离决定。

\cppfile{code/tree/hld4.cpp}

\section{树上算法}

\subsection{倍增 LCA 与换根 DP}

以父指针倍增求最近公共祖先，并在两遍树上遍历中传递子树信息。倍增层数采用源码中的固定上界。

\cppfile{code/tree/hld2.cpp}

\subsection{树链剖分与线段树}

剖分将路径拆成若干重链区间，由区间数据结构维护这些区间。

\cppfile{code/tree/hld1.cpp}

\section{常考进阶}

\subsection{树上路径交}

利用路径端点的最近公共祖先计算两条路径的公共节点数。代码中的 LCA 查询使用同一棵有根树。

\cppfile{code/tree/hld3.cpp}

\subsection{树哈希与同构}

自底向上合并孩子的哈希，用于比较树的结构。哈希判等存在碰撞的可能。

\cppfile{code/tree/hld5.cpp}
